\chapter{Learning to Move}
\label{ch:learning_to_move}

\begin{gomdraftnotice}
This chapter is under active development and contains only a structural outline
with preliminary material. Full exposition, worked examples, proofs, and exercises
will be added in future editions.
\end{gomdraftnotice}

\epigraph{%
  The first swing establishes the manifold. The ten-thousandth swing perfects the funnel.
}{}

\vspace{0.5cm}

\begin{laymansbox}{Practice Makes Funnels}{practice_makes_funnels}
When you swing a golf club for the first time, your brain is guessing how much force to use, guessing the weight of the club, and guessing how your muscles will respond. The mathematical models we've built so far are just like that first swing—a really smart guess, but still a guess. Real physics (wind, friction, fatigue) is too messy to model perfectly. This chapter is about how systems *learn*. Instead of trying to calculate the perfect math equations from scratch every time, the system tries a movement, measures how badly it missed the target, and updates its "muscle memory" for the next try. Over thousands of repetitions, the brain (or the computer) learns the true physics of the motion without ever needing the exact equations. It builds a whole library of these learned motions, letting it instantly adjust to new situations without thinking.
\end{laymansbox}

% ══════════════════════════════════════════════
\section{Iterative Learning Control (ILC)}
% ══════════════════════════════════════════════

Trajectory optimization (Chapter 6) computes a reference motion strictly from a mathematical model. However, dynamic models of humans—and even robots—are invariably flawed. Frictional forces are unmodeled, inertias are estimated, and actuators possess hidden dynamics. If you execute a computed reference $\traj^*(t)$ open-loop, the physical system will likely fail to strike the moving target due to these discrepancies.

\textbf{Iterative Learning Control (ILC)} resolves this by leveraging the repetitive nature of tasks like a golf swing. ILC exploits the premise that while models are inaccurate, \emph{the unmodeled dynamics are highly deterministic across repeated trials}.

\begin{remark}{ILC and the Orbit of Feasible Trajectories}{ilc_orbit}
  From the perspective of geometric control theory as presented by Agrachev and Sachkov \cite{AgrachevSachkov2004}, each iterative trial of ILC represents an exploration within the orbit of feasible trajectories generated by the control Lie algebra. The true plant's dynamics define an actual (unknown) orbit in state space; the nominal model defines an approximate orbit. ILC learns the discrepancy between these orbits by iteratively modulating the feedforward command. The convergence of ILC is thus a convergence toward the true orbit boundary---the set of trajectories reachable on the actual (unknown) plant. Persistent excitation conditions required for identifiability of model parameters correspond to activating the control Lie algebra sufficiently to explore all dimensions of the reachable subspace. In other words, ILC's success in learning motion requires that the control inputs, applied across trials, persistently excite the full Lie algebra of admissible directions.
\end{remark} 

If a robot under-swings by $5$ cm on the first attempt, the controller records this terminal error and updates the feedforward command $\control_{\text{ff}}$ for the \emph{second} attempt. Over repeated trials, the system learns the inverse dynamics of the true physical plant without ever explicitly identifying the model parameters. 

\begin{definition}{ILC Update Law}{ilc_law}
  Let $j$ denote the trial index. The feedforward command for trial $j+1$ is updated based on the command and the \emph{transverse error} $\bm{\xi}$ from trial $j$:
  \begin{equation}
    \control_j(t) = \control_{\text{ff}, j}(t) + \bm{K}_\perp(t) \bm{\xi}_j(t)
  \end{equation}
  \begin{equation}
    \control_{\text{ff}, j+1}(t) = \control_{\text{ff}, j}(t) + \bm{L}(t) \bm{\xi}_j(t)
  \end{equation}
  where $\bm{L}(t)$ is a learning filter. This updates the nominal trajectory to iteratively cancel out deterministic disturbances.
\end{definition}

\begin{algorithm}[htbp]
\caption{Iterative Learning Control Around a Nominal Swing}
\label{alg:ch10:ilc}
\begin{algorithmic}[1]
\Require Nominal feedforward command $\control_{\mathrm{ff},0}(t)$, stabilizing feedback $\bm{K}_\perp(t)$, learning filter $\bm{L}(t)$.
\Ensure A learned command sequence with reduced trial-to-trial error.
\For{trial $j = 0, 1, 2, \ldots$}
  \State Execute the closed-loop swing:
  \State \quad $\control_j(t) = \control_{\mathrm{ff},j}(t) + \bm{K}_\perp(t)\bm{\xi}_j(t)$.
  \State Measure the transverse tracking error $\bm{\xi}_j(t)$ and the terminal miss at impact.
  \State Filter the measured error to remove non-repeatable sensor noise.
  \State Update the next feedforward command:
  \State \quad $\control_{\mathrm{ff},j+1}(t) = \control_{\mathrm{ff},j}(t) + \bm{L}(t)\bm{\xi}_j(t)$.
  \If{the terminal error and update norm are below tolerance}
    \State Stop learning.
  \EndIf
\EndFor
\Return The learned feedforward command $\control_{\mathrm{ff}}^\star(t)$
\end{algorithmic}
\end{algorithm}

\begin{center}
\begin{tikzpicture}[scale=1.1]
  % X-axis: trials
  \draw[-{Stealth[length=2mm]}, thick, gray] (0, 0) -- (6.5, 0) node[right, font=\small] {Trial $j$};

  % Y-axis: error
  \draw[-{Stealth[length=2mm]}, thick, gray] (0, 0) -- (0, 4) node[above, font=\small] {Terminal Error};

  % Error with model mismatch (without learning)
  \draw[thick, chapblue!60, decoration={markings,
      mark=at position 0.3 with {\circle{1pt}},
      mark=at position 0.5 with {\circle{1pt}},
      mark=at position 0.7 with {\circle{1pt}},
      mark=at position 0.9 with {\circle{1pt}}
    }, postaction={decorate}]
    (1, 2.8) -- (6, 2.8) node[right, font=\small] {Open-loop (model error)};

  % ILC convergence
  \draw[very thick, trajgreen, decoration={markings,
      mark=at position 0.2 with {\circle{1.5pt}},
      mark=at position 0.35 with {\circle{1.5pt}},
      mark=at position 0.5 with {\circle{1.5pt}},
      mark=at position 0.65 with {\circle{1.5pt}},
      mark=at position 0.8 with {\circle{1.5pt}},
      mark=at position 0.95 with {\circle{1.5pt}}
    }, postaction={decorate}]
    (1, 3.5) .. controls (2, 2.9) and (3.5, 0.8) .. (6, 0.2);
  \node[trajgreen!80!black, font=\small] at (4.5, 1.5) {ILC + Feedback};

  % Annotations
  \node[gray, font=\tiny, align=center] at (3, 3.2) {Deterministic\\disturbances};
  \node[gray, font=\tiny, align=center] at (3, 0.6) {Learned\\dynamics};

  % Grid
  \draw[gray!20] (1, 0) grid (6, 4);
\end{tikzpicture}
\end{center}

% ══════════════════════════════════════════════
\section{Policy Search and Reinforcement Learning}
% ══════════════════════════════════════════════

While ILC updates the feedforward commands to zero-out deterministic tracking errors, \textbf{Policy Search} algorithms optimize the \emph{shape of the funnel} itself.

Instead of computing the transverse LQR gains $\bm{K}_\perp(t)$ from an analytical Riccati equation (Chapter 4), we parameterize the controller as a policy $\pi_{\bm{\theta}}(\state)$ and search the parameter space $\bm{\theta}$ to minimize the variance at the moving goal under realistic noise.

Reinforcement Learning (RL), specifically policy gradient methods, allows the agent to sample thousands of golf swings in simulation, continuously updating its policy parameters to maximize the expected reward (clubhead speed and squareness at impact). A heavily trained RL agent organically converges on the very principles we mathematically derived in earlier chapters:
\begin{enumerate}
    \item It abandons rigid time-tracking and autonomously discovers phase-variable mapping.
    \item It exploits the underactuated double-pendulum whipping motion to gain speed.
    \item It artificially widens its "funnel" during the transition phase, realizing that overly tight feedback control here merely amplifies motor noise and spoils the impact constraints.
\end{enumerate}

\subsection{Evolutionary Policy Search}

Not every learning architecture exposes gradients cleanly. If the policy contains contact logic, lookup tables, gait phase switches, or parameters chosen from expensive black-box simulation, \textbf{genetic} and related evolutionary methods become attractive \cite{Goldberg1989,StornPrice1997,HansenOstermeier2001}. Their role here is not to replace the geometry; it is to help search over controller families when differentiability is weak or deceptive.

\begin{remark}{Evolutionary Search as a Library Builder}{evolutionary_library_builder}
  The most natural place for evolutionary methods in this volume is the outer loop over \emph{strategy parameters}: swing tempo, release timing, phase-variable thresholds, or shape parameters of a feedback policy. Once a promising family is found, the resulting motion should still be analyzed with the trajectory-optimization and funnel tools from the previous chapters.
\end{remark}

% ══════════════════════════════════════════════
\section{Trajectory Libraries}
% ══════════════════════════════════════════════

A singular optimal trajectory $\traj^*$ and its corresponding funnel $\mathcal{F}$ are highly specific to a single objective. What if the golfer wishes to fade the ball, hit a draw, or execute a punch shot out of the rough?

Instead of re-running a massive NLP optimization for every scenario, biological systems and advanced controllers maintain \textbf{Trajectory Libraries}. A library consists of several distinct, pre-computed optimal trajectories and their corresponding local funnels.

When faced with a new task (e.g., striking the ball on an uneven lie), the system searches its library for the closest matching funnel. By interpolating between established funnels using techniques like LQR-Trees~\cite{Tedrake2009} or simple spatial blending, the system can instantly generate a stable, feasible motion without any online optimization.

\begin{algorithm}[htbp]
\caption{Trajectory-Library Retrieval and Adaptation}
\label{alg:ch10:trajectory-library}
\begin{algorithmic}[1]
\Require Task descriptor $\bm{z}$, library entries $\{(\traj_i^*, \mathcal{F}_i, \bm{z}_i)\}$, local adaptation routine.
\Ensure A feasible warm start for the current task.
\State Compute similarity scores between the requested task $\bm{z}$ and each stored descriptor $\bm{z}_i$.
\State Select the best matching entry or a small elite subset.
\If{one funnel already contains the current initial condition and task variation}
  \State Reuse that trajectory and controller directly.
\Else
  \State Blend nearby trajectories or policy parameters to form a warm start.
  \State Run a short local optimization or ILC phase to adapt the warm start to the new task.
\EndIf
\State Verify that the adapted trajectory still lies inside a certified funnel near impact.
\Return The adapted trajectory/controller pair.
\end{algorithmic}
\end{algorithm}

% ══════════════════════════════════════════════
\section{Summary}
% ══════════════════════════════════════════════

\begin{center}
\renewcommand{\arraystretch}{1.4}
\begin{tabular}{@{} l p{7.5cm} @{}}
  \toprule
  \textbf{Concept} & \textbf{Mechanism of Learning} \\
  \midrule
  Iterative Learning Control (ILC) & Feedforward update: $\control_{\text{ff}, j+1} = \control_{\text{ff}, j} + \bm{L}(t) \bm{\xi}_j$. Learns deterministic model mismatch. \\
  Policy Search & Direct optimization of feedback parameters $\bm{\theta}$ via gradient methods on sampled task performance. \\
  Evolutionary Search & Population-based tuning of policy or trajectory parameters when gradients are unreliable or unavailable. \\
  Reinforcement Learning (RL) & Learns both reference trajectories and feedback policies through trial-and-error in simulation or on real systems. \\
  Trajectory Library & Pre-computed collection of optimal trajectories and funnels for multiple task variants; enables instantaneous adaptation via interpolation. \\
  Model Identification & Learning the true plant dynamics $\bm{f}_{\text{true}}(\state, \control)$ to replace or refine the nominal model. \\
  \bottomrule
\end{tabular}
\end{center}

\vspace{0.3cm}
